\chapter{深度学习工具}

附录把运行环境写成可重复的映射：笔记本、云实例、硬件选择与贡献流程。
计算仍落在同一套张量运算上；差别在于谁提供 CPU/GPU 以及代码如何进入仓库。

\section{使用Jupyter Notebook}

笔记本是交错的叙述单元与代码单元。本地或远程内核执行单元，返回张量与打印。

\subsection{在本地编辑和运行代码}

工作目录指向本书代码后启动笔记本服务，浏览器连到本地端口。
打开 \texttt{.ipynb} 即可按单元求值。

\subsection{高级选项}

原生笔记本含大量运行期元数据，不利于版本合并。可用 Markdown 源编辑；远程机器则用端口转发。

\subsubsection{Jupyter中的Markdown文件}

贡献应改 Markdown 源而非 \texttt{.ipynb}。插件把 \texttt{.md} 当作笔记本打开，使公式与代码单元可原地运行。

\subsubsection{在远程服务器上运行Jupyter Notebook}

SSH 把远程内核端口映到本机同一或另一端口，浏览器仍访问 \texttt{localhost}。

\subsubsection{执行时间}

对每个代码单元记录墙钟时间，便于比较 $\bm{A}^{\mathsf T}\bm{B}$ 与 $\bm{A}\bm{B}$ 等运算。

\subsection{小结}

本地或经端口转发远程编辑单元。源文件用 Markdown 时，版本记录的是文本而不是输出缓存。

\subsection{练习}

核验：在 $\mathbb{R}^{1024\times 1024}$ 上 $\bm{A}^{\mathsf T}\bm{B}$ 与 $\bm{A}\bm{B}$ 哪一个更快。

\section{使用Amazon SageMaker}

当本地算力不够时，用托管笔记本实例跑 GPU 密集型单元。

\subsection{注册}

创建云账户并打开 SageMaker 控制台。建议打开计费警报，避免实例忘记停止。

\subsection{创建SageMaker实例}

选定实例类型即选定 $(\text{CPU 核数},\;\mathrm{GPU})$。
例如单卡 V100 加多核 CPU，足以覆盖本书大部分训练。

\subsection{运行和停止实例}

实例就绪后在浏览器打开笔记本。停止实例停止按计算计费；忘记停止会持续产生费用。

\subsection{更新Notebook}

在实例终端向远程仓库取更新。若有本地改动，先提交或丢弃再拉取，以免合并冲突。

\subsection{小结}

SageMaker 把“开 GPU 笔记本”变成创建–打开–停止。更新走实例上的 Git。

\subsection{练习}

核验：任选一节 GPU 代码在实例上能否跑通，以及终端是否指向笔记本目录。

\section{使用Amazon EC2实例}

相对托管笔记本，自建虚拟机通常更便宜。步骤：申请 GPU 机、装 CUDA、装框架、端口转发。

\subsection{创建和运行EC2实例}

在 EC2 控制台启动实例。下面按位置、配额、映像与登录展开。

\subsubsection{预置位置}

选邻近区域以降低控制面延迟。需确认该区域提供所需 GPU 机型。

\subsubsection{增加限制}

新账户对 GPU 机型常有配额 $0$。先申请提高限额，再启动。

\subsubsection{启动实例}

选 Ubuntu 类映像与 GPU 机型。不同代号对应不同 GPU 个数与显存，按任务选。

\subsubsection{连接到实例}

用密钥文件 SSH。密钥权限必须仅本人可读，否则握手失败。

\subsection{安装CUDA}

先更新驱动，再按发行说明安装与框架匹配的 CUDA 工具包。
安装包 URL 与文件名随版本变，以厂商当前仓库为准。

\subsection{安装库以运行代码}

按本书安装章在远程 Linux 上建隔离环境并安装 PyTorch 与辅助包。
无图形界面时，下载器用命令行取安装脚本，初始化后再加载 shell 配置。

\subsection{远程运行Jupyter笔记本}

云主机无显示器。从本地 SSH 并转发端口，使浏览器连本机端口即连上远程内核。
端口转发必须在本地命令行发起：本地端口映到远程笔记本端口。
激活环境后启动笔记本，把给出的 URL 中的端口改成本地转发端口再打开。

\subsection{关闭未使用的实例}

停止：盘仍在，可再开机，收少量存储费。
终止：盘与数据删除，不可恢复。
需要模板时先做映像再终止。

\subsection{小结}

按需启动 GPU 机；框架前先装 CUDA；用本地发起的端口转发跑笔记本。

\subsection{练习}

核验：竞价实例相对按需的价格，以及多卡机上数据并行能扩到几张卡。

\section{选择服务器和GPU}

训练受 GPU 小时约束。单机可到数张卡；再多通常用云，因供电与散热。

\subsection{选择服务器}

算力在 GPU 上，CPU 核数不必最多，但单线程频率在多卡加 Python 时仍重要。
电源按每卡峰值（可到数百瓦）留余量；机箱与风道必须容下长卡。
PCIe 通道数限制满速 GPU 个数。

\subsection{选择GPU}

加速库对其中一家的通用计算接口支持更成熟，故训练卡多选该路线。
消费级与企业级算力接近，企业级多为被动散热、更大显存与 ECC，单价高一个数量级。
实验室规模常用消费级；上百台则考虑企业级或云。新一代卡通常能效更好。低精度（如 FP16）提高吞吐，但训练需防止下溢。

\subsection{小结}

电源、通道、单线程 CPU 与散热决定能插几张卡。
大部署用云；核对机械尺寸与功耗是否匹配机箱。

\section{为本书做贡献}

改正笔误、链接与表述通过向源仓库提交拉取请求完成。持续集成会跑改动过的代码。

\subsection{提交微小更改}

在网页上打开对应 Markdown，改一句，写说明并开拉取请求。

\subsection{大量文本或代码修改}

源为 Markdown 加书籍扩展（公式、图、交叉引用）。
改代码时用笔记本跑通，提交前清空输出，由 CI 重算。
多框架章节用标记区分代码块所属框架。

\subsection{提交主要更改}

标准 Git：派生、克隆、本地改、推送、开拉取请求。

\subsubsection{安装Git}

在本机安装 Git 并注册托管账户。

\subsubsection{登录GitHub}

在浏览器打开本书仓库，派生到自己的命名空间。

\subsubsection{克隆存储库}

用派生仓库的地址在本地做副本，再按安装章建环境。
克隆地址中的用户名必须换成派生后的账户名，否则指向的不是可写副本。

\subsubsection{编辑和推送}

改文件后查看差异，提交到自己的分支并推送。

\subsubsection{提交Pull请求}

在派生仓库上对新改动开拉取请求，写清改了什么。可能被直接合并、要求修改或拒绝。
请求应小，便于审阅。

\subsection{小结}

小改动可网页编辑；大改动走派生–本地–拉取请求。避免巨型请求。

\subsection{练习}

核验：完成一次派生，并尝试用新分支提交最小改动。

\section{d2l API 文档}

辅助包把模型、数据、训练循环与公用函数收成可复用映射。实现见源文件。

\subsection{模型}

网络模块：$f_{\bm{\theta}}$ 的层、块与常见骨干，前向为张量到张量。

\subsection{数据}

数据集与加载器：$(\mathcal{D})\mapsto$ 小批量迭代器，含词表、增广与填充。

\subsection{训练}

训练循环：给定 $f_{\bm{\theta}}$、迭代器与优化器，执行前向、损失、反传与评估。

\subsection{公用}

作图、计时、设备选择与累加器等与任务无关的工具函数。
